\documentclass[12 pt, letterpaper]{hmcpset}
\usepackage{hw-preamble}
\usepackage{pdfpages}

\name{Jayden Thadani}
\class{MATH 145 --- Topics in Geometry and Topology}
\assignment{Homework 10}
\duedate{16th April 2025}

\begin{document}

Recall that $E = \{ \O \}$ denotes the empty simplicial complex. Recall also that $C$ and $S$ denote the cone and suspension operations.

\begin{problem}[52.]
	\begin{enumerate}
		\item Sketch the simplicial complexes $CE$, $CCE$,  $C C C E$, $C C C C E$.

		\item Determine the dimension of $C C C C C E$ and count the top-dimensional simplices.
	\end{enumerate}
\end{problem}

\begin{solution}
	\begin{enumerate}
		\item We draw the simplicial complexes with points as dots, edges as lines, triangles with squiggles, and tetrahedra shaded in.
			\begin{figure}[H]
				\centering
				\includegraphics[width = 0.85 \textwidth]{figures/P52a cones.pdf}
				\caption{In order, $CE$, $CCE$, $C C C E$, and $C C C C E$.}
			\end{figure}

		\item
			\[
				\boxed{
					\begin{split}
						\dim C C C C C E = 4 \\
						\#  C C C C C E^{(4)} = 1.
					\end{split}
				}
			\]

			We claim that any $C^{k} E$ has dimension $k - 1$ and contains exactly one $(k - 1)$-dimensional simplex. $C^{0} E = E$ satisfies the base case --- it contains exactly one $(-1)$-dimensional simplex $\O$. 

			Suppose $C^{k - 1} E$ has dimension $k - 2$ and contains exactly one $(k - 2)$-simplex $\sigma$. Then $C^{k}$ has dimension at most $k - 1$ since it only contains one more vertex than $C^{k - 1}$. Since the simplices of $C^{k} E = C C^{k - 1} E$ are just $\{ * \} \cup \tau$ for $\tau \in C^{k - 1} E$ and there is only one $(k - 2)$-simplex in $C^{k - 1} E$, there is at most one $(k - 1)$-simplex in $C^{k} E$. The simplex $\{ * \} \cup \sigma$ where $\sigma$ is the unique $(k - 1)$-simplex in $C^{k} E$.
	\end{enumerate}
\end{solution}

\pagebreak

\begin{problem}[53.]
	\begin{enumerate}
		\item Sketch the simplicial complexes $SE$, $SSE$, $SSSE$.

		\item Determine the dimension of $SSSSE$ and count the top-dimensional simplices.
	\end{enumerate}
\end{problem}

\begin{solution}
	\begin{enumerate}
		\item We draw the simplicial complexes with the same scheme as previously
			\begin{figure}[H]
				\centering
				\includegraphics[width = 0.85 \textwidth]{figures/P53a suspensions.pdf}
				\caption{In order, $SE$, $SSE$, and $SSSE$}
			\end{figure}


		\item
			\[
				\boxed{
					\begin{split}
						\dim SSSSE & = 3 \\
						\# SSSSE^{(3)} & = 16.
					\end{split}
				}
			\]
			We claim that any $S^{k} E$ (for $k$ a positive integer) has dimension $k - 1$ and contains exactly $2^{k}$ top-dimensional $(k - 1)$-simplices. $S^{1} E$ satisfies the base case --- it contains exactly $2^{1}$ $0$-simplices --- the two points added by the $S$ operation.

			Suppose $S^{k - 1} E$ has dimension $k - 2$ and contains exactly $2^{k - 1}$ top-dimensional $(k - 2)$-simplices $\sigma_{j}$. Then since $S^{k} E = S S^{k - 1} E$, each $S^{k} E$ contains two $(k - 1)$-simplices for each $\sigma_{j}$ --- $\tau_{j, 1} = \sigma_{j} \cup \{ *_{1} \}$ and $\tau_{j, 2} = \sigma_{j} \cup \{ *_{2} \}$. Thus, $S^{k} E$ has $2 2^{k - 1} = 2^{k}$ many $(k - 1)$-simplices.
	\end{enumerate}
\end{solution}

\pagebreak


\begin{problem}[54.]
	Given a poset $P$, we define two related simplicial complexes as follows. In both cases, $P$ is the vertex set. The \emph{simplicial complex of chains} $\operatorname{Ch} P$ is defined
	\[
		\sigma = \{ x_{0}, x_{1}, \dots, x_{n} \} \in \operatorname{Ch} P
	\]
	if and only if every pair of elements $x_{i}, x_{j}$ is comparable.

	The \emph{simplicial compelx of antichains} $\operatorname{Anti} P$ is defined
	\[
		\sigma = \{ x_{0}, x_{1}, \dots, x_{n} \} \in \operatorname{Anti} P
	\]
	if and only if no pair of distinct elements $x_{i}, x_{j}$ is comparable.

	Let $P$ be the poset defined by the following \emph{Hasse diagram}:

	In other words, $a$ is greater than everything else, $e$ is less than everything else, $b, c, d$ are pairwise incomparable with each other.

	\begin{enumerate}
		\item Describe and sketch $\operatorname{Ch} P$.

		\item Describe and sketch $\operatorname{Anti} P$.
	\end{enumerate}
\end{problem}

\begin{solution}
	Note that from now we draw triangles shaded in rather than with squiggles since we don't have to draw any more tetrahedra.
\begin{enumerate}
	\item It suffices to look at the maximal chains of the poset --- $\operatorname{Ch} P$ is the complex generated by the triangles corresponding to chains $eca, eba, eda$.

		\begin{figure}[H]
			\centering
			\includegraphics[width = 0.5 \textwidth]{figures/P54a chain.pdf}
			\caption{$\operatorname{Ch} P$ drawn with triangles shaded in}
		\end{figure}

	\item It suffices to look at the maximal sets so that all elements are pairwise incomparable. These are $\{ b, c, d \}, \{ a \}, \{ e \}$. Note that points are counted as anti-chains since they are vacuously pairwise incomparable. Thus, $\operatorname{Anti} P$ consists of the triangle $bcd$ and points $a$ and $e$
		\begin{figure}[H]
			\centering
			\includegraphics[width = 0.5 \textwidth]{figures/P54b anti.pdf}
			\caption{$\operatorname{Anti} P$ drawn with triangles shaded in}
		\end{figure}
\end{enumerate}
\end{solution}

\pagebreak

\begin{problem}[55.]
	\begin{enumerate}
		\item Draw a regular hexagon whose vertices lie on the unit circle in the plane, and calculate the distances which occur between pairs of these vertices.

		\item For all values of $R \ge 0$, identify the Vietoris-Rips complex $\operatorname{VR}(X, R)$ where $X$ is the set of vertices of a regular hexagon.
	\end{enumerate}
\end{problem}

\begin{solution}
	\begin{enumerate}
		\item The side distance is $1$ since together with the radii to the endpoints of each side, we have an equilateral triangle. The distance between a point on the hexagon and its opposite is $2$ since it is a diameter of the unit circle.

			We can get the length of the other diagonal of the hexagon by trigonometry. Half of the diagonal forms a right angled triangle with the radius of the circle as the hypotenuse and the angle opposite $\pi / 3$. Thus, half the diagonal has length $\sin{(\pi / 3)} = \sqrt{3} / 2$, and the diagonal has double that length --- it has length $\sqrt{3}$. By symmetry these are all the distances.

			\begin{figure}[H]
				\centering
				\includegraphics[width = 0.5 \textwidth]{figures/P55a lengths.pdf}
				\caption{All the distances between vertices of the hexagon}
			\end{figure}

		\item For $R < 1$ no points are connected.
			\begin{figure}[H]
				\centering
				\includegraphics[width = 0.5 \textwidth]{figures/P55b vertices.pdf}
				\caption{$\operatorname{VR}(X, R)$ for $R < 1$}
			\end{figure}

			For $1 \le R < \sqrt{3}$, we get the boundary of the hexagon.
			\begin{figure}[H]
				\centering
				\includegraphics[width = 0.5 \textwidth]{figures/P55b boundary.pdf}
				\caption{$\operatorname{VR}(X, R)$ for $1 \le R < \sqrt{3}$}
			\end{figure}

			For $\sqrt{3} \le R < 2$, we get a filled in boundary of the hexagon that is homtopy equivalent to the boundary of the hexagon.
			\begin{figure}[H]
				\centering
				\includegraphics[width = 0.5 \textwidth]{figures/P55b filled.pdf}
				\caption{$\operatorname{VR}(X, R)$ for $\sqrt{3} \le R < 2$}
			\end{figure}

			For $R \ge 2$ we get the whole hexagon.
			\begin{figure}[H]
				\centering
				\includegraphics[width = 0.5 \textwidth]{figures/P55b all.pdf}
				\caption{$\operatorname{VR}(X, R)$ for $R \ge 2$}
			\end{figure}
	\end{enumerate}
\end{solution}

\pagebreak

\begin{problem}[56.]
	\begin{enumerate}
		\item Exhibit a simplicial complex $P$ with $\mathrm{b}_{0}(P) = 4$ and $\chi(P) = 7$.

		\item Exhibit a simplicial complex $P$ with $\mathrm{b}_{0}(P) = 7$ and $\chi(P) = 4$. 
	\end{enumerate}
\end{problem}

\begin{solution}
	Note that the Euler characteristic of the disjoint union of simplicial complexes is just the sum of their Euler characteristics. Note also that a point has Euler characteristic $1$, the $2$-sphere --- boundary of the $3$-ball has Euler characteristic
	\[
		\chi(S^{2}) = 4 - 6 + 4 = 2
	\]
	and the homotopy circle ($S^{1} \ = C(\{ 1, 2, 12 \}, \{ 1, 2 \})$) has Euler characteristic
	\[
		\chi(S^{1}) = 3 - 3 = 0.
	\]
	These are all connected simplicial complexes. Thus, by taking disjoint unions of them, we can create simplicial complexes with desired zeroth Betti number and Euler characteristic.
	\begin{enumerate}
		\item The disjoint union of three $S^{2}$s and a point has Euler characteristic
			\[
				3  \chi(S^{2}) + \chi(\{ * \}) = 7.
			\]
			and is the disjoint union of four connected simplicial complexes, so has zeroth Betti number $4$.
			\begin{figure}[H]
				\centering
				\includegraphics[width = 0.5 \textwidth]{figures/P56a b4chi7.pdf}
				\caption{Three ``spheres'' and a point}
			\end{figure}

		\item The disjoint union of three $S^{1}$s and four points has Euler characteristic
			\[
				3 \chi(\mathbb{T}) + 4 \chi(\{ * \}) = 4
			\]
			and is the disjoint union of seven connected simplicial complexes, so has zeroth Betti number $7$.
			\begin{figure}[H]
				\centering
				\includegraphics[width = 0.5 \textwidth]{figures/P56b b7chi4.pdf}
				\caption{Three ``circles'' and four points}
			\end{figure}
	\end{enumerate}
\end{solution}

\pagebreak

\begin{problem}[57.]
	Let $X$ be a simplicial complex $X$. Removing the empty simplex, we consider $X \setminus \{ \O \}$ as a poset with $\sigma \le \tau$ if $\sigma$ is a face of $\tau$. We define the \emph{derived complex of $X$}
	\[
		 X^{(1)} = \operatorname{Ch} (X \setminus \{ \O \})
	\]
	to be the simplicial complex of chains of $X \setminus \{ \O \}$.

	In each part, sketch the derived complex of the given complex. Try to find a good way to draw it that reveals what is actually going on with this construction.

	\begin{enumerate}
		\item $X = B_{0}$.

		\item $X = B_{1}$.

		\item $X = \partial B_{2}$.

		\item $X = B_{2}$.

		\item $X = $
			\begin{figure}[H]
				\centering
				\includegraphics[width = 0.45 \textwidth]{figures/P57e X.pdf}
			\end{figure}
	\end{enumerate}
\end{problem}

\begin{solution}
	Taking the derived complex looks a lot like barycentric division.
	\begin{enumerate}
		\item Taking the derived complex of a point leaves it unchanged
			\begin{figure}[H]
				\centering
				\includegraphics[width = 0.45 \textwidth]{figures/P57a B0.pdf} \includegraphics[width = 0.45 \textwidth]{figures/P57a B0.pdf}
				\caption{$B_{0}$ and its derived complex}
			\end{figure}

		\item Taking the derived complex of a line turns it into two lines
			\begin{figure}[H]
				\centering
				\includegraphics[width = 0.45 \textwidth]{figures/P57b B1.pdf} \includegraphics[width = 0.45 \textwidth]{figures/P57b B1der.pdf}
				\caption{$B_{1}$ and its derived complex}
			\end{figure}

		\item We can use this to see that taking the derived complex of the boundary of a triangle turns it into the boundary of a hexagon
			\begin{figure}[H]
				\centering
				\includegraphics[width = 0.45 \textwidth]{figures/P57c delB2.pdf} \includegraphics[width = 0.45 \textwidth]{figures/P57c delB2der.pdf}
				\caption{$\partial B_{2}$ and its derived complex}
			\end{figure}

		\item Thus, the derived complex of the triangle is a hexagon with triangles in it
			\begin{figure}[H]
				\centering
				\includegraphics[width = 0.45 \textwidth]{figures/P57d B2.pdf} \includegraphics[width = 0.45 \textwidth]{figures/P57d B2der.pdf}
				\caption{$B_{2}$ and its derived complex}
			\end{figure}

		\item Combining the derived complex of $B_{2}$ and $\partial B_{2}$ gives the following derivation
			\begin{figure}[H]
				\centering
				\includegraphics[width = 0.45 \textwidth]{figures/P57e X.pdf} \includegraphics[width = 0.45 \textwidth]{figures/P57e Xder.pdf}
				\caption{$X$ and its derived complex}
			\end{figure}
	\end{enumerate}
\end{solution}

\pagebreak

\begin{problem}[58.]
	Consider the Whitehead equivalence relation, generated by elementary collapses. Two of the four complexes drawn below are Whitehead equivalent to each other, as are the other two.

	For each equivalent pair, demonstrate the Whitehead equivalence by exhibiting a sequence of elementary collapses and expansions that lead from one to the other.

	Note that in the second Whitehead sequence we use the edge-splitting lemma.
\end{problem}

\begin{solution}
	The Whitehead sequences are included below. Expansions and collapses are indicated in colour, with collapses indicated by drawing the removed face as dotted (without drawing anything for the removed coface).
	\includepdf[pages=-]{figures/P58a whitehead.pdf}
	\includepdf[pages=-]{figures/P58b whitehead.pdf}
\end{solution}

\end{document}
